% 第 18 章　声音与信息
\chapter{声音与信息}

\step{反常识开场}

\begin{mainline}
找一首你非常熟悉的歌，用手机外放，闭上眼睛听十秒钟。

想一想你刚才听到了什么：鼓、贝斯、吉他、键盘，也许还有和声，几十件乐器层层叠叠。可是你的手机只有一个扬声器，扬声器里只有一片薄薄的膜。在任何一个瞬间，这片膜只在一个位置，只朝一个方向运动——要么往外推一点，要么往里收一点。

问题来了：一片只会“前一点、后一点”的膜，怎么装得下一整个乐队？

更离奇的是，你不光能听出“有鼓有吉他”，还能听出是木吉他还是电吉他，能听出唱歌的人是谁，甚至能听出他今天心情怎么样。这些信息全都只能从那一片膜的一来一回里挤出来。

同样的怪事藏在电话里。对方的声音从听筒出来时，可能已经旅行了上千公里，可你听第一句话就能认出是谁。一路上被搬来搬去的到底是什么？为什么一阵简简单单的“抖动”，能装下姓名、语句和情绪？

这一章要把声音拆到底：它是什么东西在传播，它为什么能携带信息，以及——它最多能携带多少。走完这一章你会发现，“一片膜装下一个乐队”不但不是魔术，反而是通信技术每天都在做的事。
\end{mainline}

\step{先猜一猜}

\begin{mainline}
往下读之前，先押三个注。

第一注。把一只正在响的闹钟放进玻璃罩里，用抽气机慢慢抽走罩里的空气。你猜会听到什么？声音会一直保持原样，还是越来越小？如果抽得足够干净，最后会怎样？

第二注。你和朋友隔着五十米。你先用最大的嗓门喊他的名字，再用平常聊天的音量说一遍。两次声音到达他耳朵所用的时间，哪一次更短？还是一样长？

第三注。录音机（或者手机）把一首歌存下来，你觉得它存的是一份“乐器名单”——这里有一面鼓、那里有一把吉他，到播放时再把它们合成出来？还是存的是别的东西？

把答案写下来。这一章走完，你自己给这三个注评分——提前声明：至少有一注，答案会和多数人的直觉相反。
\end{mainline}

\step{动手实验}

\begin{experiment}[棉线电话和桌边的尺子]
找两个纸杯（或两个空罐头盒）、一根五米以上的棉线。在杯底正中各戳一个小孔，把棉线两端分别穿进去，在线头系上一根牙签，让它卡在杯底。你和朋友各拿一个杯子走开，把棉线拉直绷紧，然后一个人对着杯口小声说话，另一个人把杯口罩在耳朵上听。

你会听到：声音清楚得出奇，比隔着空气直接喊清楚得多。

接着做三个破坏性的检查。第一，让棉线松垮下来，声音立刻变得含糊甚至消失。第二，在线绷紧的时候，请第三个人用手指轻轻捏住棉线中点——声音马上哑掉。第三，把棉线换成同样长的毛线，再换成橡皮筋，比较效果。

然后换一个小实验。把一把钢尺压在桌边，让一截伸出桌面，按住桌面上的部分，另一只手拨一下悬空端。钢尺振动，发出一个音。把伸出部分缩短一半，再拨一次——音调明显变高了。

如果还有兴致，再找三个一样的玻璃杯，装入不同深度的水，用筷子挨个轻敲杯壁。水装得越满，音调越低。这次振动的不只是玻璃，而是杯子连水一起的整体，“个子”越大，振动越慢。

请记下三个观察：绷紧是传声的必要条件；任何东西碰到传声的路径，声音就会被“截走”；振动部分越短、越轻，音调越高。这一章后面要逐条解释它们。
\end{experiment}

\begin{mainline}
这两个实验加在一起，已经否定了不少想当然。杯子之间的棉线没有输送任何“东西”过去——线还是那根线，一点没少走。可线一松、一被捏住，通道就断了。而钢尺和水杯告诉你：音调高低和振动的快慢有直接关系，振动的“个子”越小，振动越快。

还有一个不用道具的检查，现在就做：把耳朵贴在桌面上，用指甲在桌面的另一头轻轻刮一下。声音又大又清楚，仿佛刮在你的耳膜上；抬起耳朵，在空气里听同一个动作，几乎听不见。固体传声的效率，比空气高得多——印第安人把耳朵贴在地面上听远处的马蹄声，用的是同一个道理。
\end{mainline}

\step{旧模型遇到麻烦}

\begin{mainline}
每个人脑子里都有一个默认的声音模型，可以叫“包裹模型”：发声体把声音打包好，交给空气，空气像快递员一样把它送到你的耳朵。这个模型日常生活里很好用，但只要认真追问，它连出三个洋相。

第一个洋相：速度对不上。声音在空气里每秒跑约三百四十米。如果真是空气把声音“送”过去的，那就该有一股每秒三百四十米的气流从你的嘴边刮到听者的脸上——这是十二级台风风速的十倍。实际上说话时连对方额前的头发都吹不乱。包裹到了，快递员却没动。

第二个洋相：双向对不上。你和朋友面对面同时说话，两股声音迎面穿过同一段空气，互不干扰，各自到达。如果声音是“一股气流”，两股对撞的气流只会搅成一团。可声音没有搅成一团——你们各自听清了对方的话，不然吵架就吵不起来了。

第三个洋相最致命：真空。把响着的闹钟封进玻璃罩，抽走空气，铃声会越来越弱，最后几乎消失——隔着玻璃你明明能看见铃锤还在敲打。看得见的振动传得出光，却传不出声。抽走的是空气，断掉的是声音。这说明空气不是“快递员”那么简单，它是这条路上不可或缺的“路面”本身：没有它，声音连一步都迈不出去。

顺带还有个小洋相：山谷里的回声。包裹送到就该签收消失，为什么声音撞上山壁会原路弹回来，有时还弹回来好几次？“包裹模型”对此一言不发。

三个洋相加起来指向同一个结论：在传播的从来不是“一团东西”，而是一种“状态”——这正是第 17 章里波的样子。体育场人浪里没有观众跑圈，传的是“站起来”这个状态；绳波里没有绳子远行，传的是“鼓包”这个形状。声音也应该是一种波。可水波的波形看得见，绳波的形状摸得着，声音的“波形”藏在哪里？要回答它，得先发明几个新词。
\end{mainline}

\step{发明新概念}

\begin{mainline}
请想象一根很长的软弹簧铺在地板上。你在这一头沿弹簧方向猛地一推一拉，会看到一圈“挤紧”沿着弹簧跑向那头，后面跟着一圈“松开”。注意：没有任何一圈弹簧从你这头跑到那头，跑过去的只是“挤紧”和“松开”这两个状态，每一圈弹簧都只在原地附近前后挪动。这种振动方向与传播方向一致的波，叫\textbf{纵波}；第 17 章里绳子上的鼓包是振动方向与传播方向垂直的横波。

现在把弹簧撤掉，换上空气。空气分子之间当然没有真的小弹簧，但效果上很像：你把一团空气压缩，它的压强升高，会反推回来，把邻居也压紧；邻居再压紧它的邻居。“挤一下、松一下”就这样一棒接一棒传下去。你说话时，声带（和嘴唇、舌头）做的就是这件事：周期性地把面前的空气挤紧、放松。于是空气里出现一列向前跑的“密”和“疏”，压强一高一低地交替。这就是\textbf{声音：介质中的疏密波，一种纵波}。被传来传去的那种“有弹性的东西”，叫\textbf{介质}——空气、水、钢轨、棉线都可以是介质，真空里没有介质，所以闹铃哑了。

图 18-1 画出了某一瞬间空气分子的分布：有的地方挤成一团（密部，压强高），有的地方稀稀拉拉（疏部，压强低）。整个图案随时间向前平移，单个分子却只在原地附近前后抖动。

这套图像对别的介质同样成立。游泳时把耳朵埋进水里，岸上原本模糊的人声会突然变得清楚——水也是介质，而且分子之间推得比空气干脆得多，接力更快。潜水的人还能听见远处轮船螺旋桨的嗡鸣，而水面上的人毫无察觉：声音在水下传得又快又远。海洋里甚至存在天然的“声道”：在几百米深的水层中，声速随深度的变化会把声波弯回这层水里，让它沿着这层水走几百上千公里而不散掉，鲸的歌声就是这样横跨大洋的。地震更是极致的例子：大地本身就是介质，一次地震掀起的疏密波和摇晃波可以绕地球跑好几圈，被半个地球外的仪器记录下来。
\end{mainline}

\begin{figure}[htbp]
\centering
\begin{tikzpicture}[x=1cm,y=1cm]
  \foreach \x in {0,0.26,0.48,0.66,1.30,2.00,2.70,2.96,3.18,3.36,4.00,4.70,5.40,5.66,5.88,6.06,6.70,7.40}
    \fill (\x,0) circle (1.4pt);
  \draw[double distance=1.5pt,-{Stealth}] (0.05,0.45) -- (0.7,0.45);
  \node[above] at (0.38,0.5) {密部};
  \draw[double distance=1.5pt,-{Stealth}] (1.4,0.45) -- (2.6,0.45);
  \node[above] at (2.0,0.5) {疏部};
  \draw[-{Stealth},thick] (6.6,-0.45) -- (7.6,-0.45) node[midway,below]{传播方向};
  \node[below] at (3.2,-0.6) {每个分子只在原地附近前后抖动，向前跑的是疏密图案};
\end{tikzpicture}
\caption{声音是空气里的疏密波（纵波）}
\end{figure}

\begin{mainline}
这里要交代一个比喻的边界。“空气像弹簧”这个说法，像的地方在于：空气被压缩时会反推回来，这份“弹性”让挤压状态可以接力传递。不像的地方在于：分子之间并没有连着的小弹簧，弹性是无数次分子碰撞统计出来的效果；而且这根“空气弹簧”只能贴着传，分子够不到的地方——比如真空——它就无计可施。

有了“波”这个概念，听觉里的三个老朋友可以各就各位了。

\textbf{音调}对应波的频率：每秒挤多少次。频率高，音调高；频率低，音调低。钢尺伸出的部分越短，振动越快，频率越高——实验里的观察对上了。

\textbf{响度}对应波的振幅：每次挤压有多狠，压强起伏有多大。使劲喊和轻轻说，挤空气的节奏可以一样，但挤的狠劲完全不同。

\textbf{音色}是最有意思的一个。同样一个音，音叉、钢琴、小提琴、人唱出来，你绝不会搞混。频率一样、响度一样，区别在哪里？区别在于波形的“花样”：真实乐器发出的从来不是一条干干净净的正弦曲线，而是疙疙瘩瘩的复杂曲线。十九世纪的法国数学家傅里叶证明了一件了不起的事：\textbf{任何复杂的周期波形，都可以拆成一组简单正弦波的叠加}——一个“主音”（基频）加上它两倍、三倍、四倍……频率的一串“陪音”（泛音），各强各弱。乐器之间的区别，就是这串陪音的配方不同。

你自己就是一台随时换配方的发声器。发“啊”和发“咿”，声带振动的节奏可以完全一样，只是你的舌头和嘴改变了口腔的形状，从而改变了陪音的配方。这就是为什么压着嗓子说悄悄话，别人照样听清你说了什么：基频可以弱到几乎没有，配方还在。

这个比喻也要划清边界。说音色像“配方”，像的地方在于：整体风味由各种成分的比例决定。不像的地方在于：配方里的成分会互相反应，而声音里的各个频率成分在空气里互不干扰——它们只是叠加，各走各的。正是这种“互不干扰”，让一片膜能装下一个乐队，这是第九节要揭开的谜底。
\end{mainline}

\step{一句话模型}

\begin{oneline}
声音不是空气跑过去，而是“挤一下、松一下”这个状态在介质里接力：音调是接力的节奏快慢（频率），响度是每次挤的轻重（振幅），音色是节奏里藏的花样（一堆简单频率的配方）。听懂声音，就是读出这串起伏里携带的信息。
\end{oneline}

\step{数学翻译器}

\begin{mainline}
第 17 章已经把波的基本关系摆在了桌上，声音直接继承：

\[
  v = f\lambda
\]

左边是波速 \(v\)：疏密图案每秒向前推进多少米。右边是频率 \(f\)（每秒挤几次）乘以波长 \(\lambda\)（一次完整的“密—疏”占多少米）。乘起来正好是“图案每秒扫过多少米”，天经地义。

先用特殊情况检查它。若 \(f\) 趋近于零，波长趋于无穷大——那已经不是声音，而是气压极其缓慢的整体升降，比如天气系统过境。若波长极短，频率就极高：在空气里（\(v\) 约每秒三百四十米），波长一厘米对应三万四千赫兹，早已超出人耳上限。再检查依赖关系：在温度给定的空气里，\(v\) 近似是个常数，所以频率翻一倍，波长必然缩一半——高音的波长就是比低音短，这件事马上有用。

人耳能听到的频率大约从二十赫兹到两万赫兹。用 \(v=f\lambda\) 折算成波长：二十赫兹对应十七米，两万赫兹对应一点七厘米。这解释了一个日常经验：隔壁房间说话，你听不清字句，却能听见嗡嗡的低音。低音波长好几米，比门缝、墙壁拐角大得多，容易绕过去（第 17 章的衍射）；高音波长短，遇到障碍就被挡下或吸收。所以楼上装修的电钻声，传到你家只剩“咚、咚”的低频部分。

再看一组真实数据。声音在十五摄氏度的空气里约每秒三百四十米，在淡水里约每秒一千四百八十米，在钢铁里约每秒五千米。规律很朴素：介质越“硬”、分子间推得越干脆，挤压状态接力就越快。这里有一个容易猜错的反例：钢比空气致密八千倍，直觉会赌声音在钢里更慢，实际上快了十几倍。决定接力快慢的是两股力量的较量——弹性让被挤的分子立刻反推回去，惯性让分子懒得动。钢的分子虽然重得多，但分子间的“推”干脆了太多太多，弹性大获全胜。所以“越密越慢”是错的，“越硬越快”才对。

现在你明白了耳朵贴铁轨的用处：一公里外的火车，声音沿钢轨约零点二秒就到，沿空气要将近三秒——贴轨听声能提前两分多钟的量级是夸张，提前两三秒是实打实的，足够预警。

空气中的声速也不是铁板一块：温度每升高一摄氏度，声速约增加零点六米每秒。热天里分子跑得更欢，挤压的接力自然传得更快。这意味着“数秒测雷”在冬天会略微估近——冷空气里声音走得慢，你以为雷电在三公里外，其实更远一些。

打雷是另一把随身尺子。闪电的光每秒三十万公里，走三公里只用十万分之一秒，可以当作“立刻到达”；雷声每秒约三百四十米，走一公里约三秒。所以看到闪电后默数秒数，除以三，就是雷电距离你的公里数。数到九秒还没听见雷，说明闪电在三公里外，这场雨暂时浇不到你。

回声是第三把尺子。对着远处的大楼喊一嗓子，约零点六秒后听到回声。声音走了往返两趟，单程零点三秒，乘以每秒三百四十米，楼就在大约一百米外。这个“发声—计时—折半—乘声速”的流程，是本章后面所有回声技术的原型：蝙蝠在黑暗中用它追捕飞蛾，轮船用它测量水深，汽车用它探测后方的障碍物。差别只在于，它们用的往往是人耳听不到的超声，计时也从你的脑子换成了电子电路。

接下来是两个本章特有的公式。第一个叫\textbf{拍频}。两把吉他弦，一把调成四百四十赫兹，另一把没调准，是四百四十二赫兹。两列波叠加后，压强起伏会像有两个节拍器在打架：每秒有两回，两列波的“密”刚好对齐，声音变大；有两回刚好错开，声音变小。于是你听到音量以每秒两次的频率“嗡——嗡——”地涨落。这个涨落频率叫拍频：

\[
  f_{\text{拍}} = \lvert f_1 - f_2 \rvert
\]

检查特殊情况：两根弦完全一样准，\(f_1=f_2\)，拍频为零——没有涨落，听到的是一个平稳的音，调音师就靠听这个“嗡”消失来校音。若两频率差很大，比如差两千赫兹，拍频高得耳朵根本跟不上，于是就听成两个独立的音——公式自动告诉你“拍”这个概念什么时候失效。

第二个公式描述救护车的警笛为什么“迎面尖叫、擦肩而过时陡然降低、远去时低沉”。这叫\textbf{多普勒效应}。设声速为 \(v\)，警笛原本的频率为 \(f\)，警笛车以速率 \(u\) 朝静止的你开来。由于声源在“追赶”自己发出的波，前方的波被挤得更密，波长变短，你听到的频率变高：

\[
  f' = \frac{v}{v-u}\,f
  \qquad\text{（声源接近；远离时把 \(u\) 换成 \(-u\)）}
\]

先做三项检查。其一，\(u=0\) 时 \(f'=f\)，车不动音调不变，废话必须成立——成立。其二，方向反转：远离时 \(u\) 变负，分母变成 \(v+u\)，频率变低——接近升高、远离降低，不对称的两个方向被同一个式子收编。其三，趋于极端：当 \(u\) 逼近 \(v\)，分母逼近零，\(f'\) 冲向无穷大。公式在尖叫求救——它在说“声源追上自己的波”这件事超出了本模型的管辖范围。事实上飞机超音速时，所有波前在机头方向叠成一堵压强的“墙”，形成\textbf{冲击波}，传到地面就是音爆。这不是“更尖的音调”，而是一种全新的东西，见“模型的边界”一节。
\end{mainline}

\begin{figure}[htbp]
\centering
\begin{tikzpicture}[x=1cm,y=1cm]
  \draw (0,0) circle (1.5);
  \draw (0.45,0) circle (1.15);
  \draw (0.9,0) circle (0.8);
  \draw (1.3,0) circle (0.45);
  \fill (1.85,0) circle (2.2pt);
  \draw[-{Stealth},thick] (1.0,-1.05) -- (2.2,-1.05) node[midway,below]{声源运动方向};
  \node at (-1.9,1.3) {波前在前进方向被挤密};
  \draw[-{Stealth}] (-1.35,1.15) -- (-0.55,0.55);
\end{tikzpicture}
\caption{运动声源发出的波前：前方挤密（频率升高），后方拉稀（频率降低）}
\end{figure}

\begin{mathbridge}
响度的定量描述藏着一个意外：人耳对“多重”的感受不是线性的，而是对数的。把声音的强度（每秒通过一平方米的能量）增大十倍，你只觉得“响了一格”；再增大十倍，又只“响一格”。所以工程上用对数刻度——分贝：

\[
  L = 10\lg\frac{I}{I_0},
  \qquad I_0 = 10^{-12}\ \text{瓦每平方米（人耳刚能察觉的强度）}
\]

算一笔账：普通谈话约六十分贝，喷气飞机起飞时约一百二十分贝。差六十分贝，强度差 \(10^6\) 倍——一百万倍。你的耳朵把一个跨度上千万倍的物理量，压缩进“安静到震耳”的十几格感受里。这是生物进化给出的数据压缩方案。

傅里叶分解写出来是这样：任何以 \(T\) 为周期的波形 \(s(t)\)，都可以写成

\[
  s(t) = a_0 + \sum_{n=1}^{\infty}\bigl[a_n\cos(n\omega t) + b_n\sin(n\omega t)\bigr],
  \qquad \omega = \frac{2\pi}{T}
\]

其中系数 \(a_n, b_n\) 记录频率为基频 \(n\) 倍的成分各有多强。把每个频率的强度画成竖条，就得到这张波形的\textbf{频谱}——声音的“成分表”。图 18-3 左边是一条复杂曲线，由一个基频和它的三倍频叠加而成；右边是它的频谱，只有两根竖条。波形在时间上“缠成一团”的信息，在频谱上一目了然。
\end{mathbridge}

\begin{figure}[htbp]
\centering
\begin{tikzpicture}[x=0.7cm,y=0.5cm]
  \draw[->] (0,3.6) -- (6.6,3.6);
  \draw[blue!60!black,domain=0:6.283,samples=80] plot (\x,{3.6+1.0*sin(\x r)});
  \node[left] at (0,4.6) {基频};
  \draw[->] (0,1.2) -- (6.6,1.2);
  \draw[blue!60!black,domain=0:6.283,samples=80] plot (\x,{1.2+0.45*sin(3*\x r)});
  \node[left] at (0,2.0) {三倍频};
  \draw[->] (0,-1.6) -- (6.6,-1.6);
  \draw[red!60!black,thick,domain=0:6.283,samples=120] plot (\x,{-1.6+1.0*sin(\x r)+0.45*sin(3*\x r)});
  \node[left] at (0,-0.6) {叠加};
  \begin{scope}[xshift=8.4cm]
    \draw[->] (0,-1.9) -- (0,4.9);
    \draw[->] (0,-1.9) -- (3.6,-1.9);
    \draw[red!60!black,very thick] (0.8,-1.9) -- (0.8,3.1);
    \draw[blue!60!black,very thick] (2.4,-1.9) -- (2.4,0.35);
    \node[below] at (0.8,-2.0) {\(f\)};
    \node[below] at (2.4,-2.0) {\(3f\)};
    \node[above,rotate=90] at (-0.25,1.5) {强度};
  \end{scope}
\end{tikzpicture}
\caption{复杂波形 = 简单正弦波的叠加（左）；频谱是它的成分表（右）}
\end{figure}

\begin{challenge}
冲击波的形状可以直接算出来。声源以速率 \(u>v\) 运动时，不同时刻发出的波前的包络围成一个圆锥（马赫锥），半顶角 \(\theta\) 满足

\[
  \sin\theta = \frac{v}{u} = \frac{1}{M},
\]

其中 \(M=u/v\) 叫马赫数。检查：\(M=1\) 时 \(\theta=90^\circ\)，锥面退化成一堵平面墙，对应“刚追上自己的波”的临界；\(M\) 越大锥越尖。快艇的尾迹、子弹的激波、超音速飞机的音爆，全是同一个几何。

还有一个更深的事实：一段声音持续的时间越短，它的频谱就摊得越宽。数学上可以证明，持续时间 \(\Delta t\) 与频率展宽 \(\Delta f\) 的乘积有一个下限，量级为

\[
  \Delta f\,\Delta t \gtrsim 1 .
\]

所以一声短促的“哒”没有明确音调，而长鸣的音叉音调很准。请注意边界：这纯粹是傅里叶数学对任何波都成立的结论，和量子力学的不确定性原理共享同一个数学骨架，但物理来源不同——那里不是“波”造成的，而是量子状态本身的结构。这个伏笔会在量子部分回收。
\end{challenge}

\step{模型的边界}

\begin{mainline}
声音模型好用，但每一条都有管辖范围。

第一，叠加原理只在小起伏时成立。日常说话、音乐，空气的压强起伏不到大气压的万分之一，此时“两列波互不干扰地相加”非常精确。但爆炸近处，压强起伏可以和大气压本身相比，空气的“弹性”不再是均匀的，波形会边走边变形，声速也不再是常数——冲击波就是这个非线性世界的居民。线性声波理论在那里整体失效。

第二，“声速与频率无关”只是很好的近似。空气里它好得惊人，所以乐队前排和后排听到的音高一致。但在固体和许多液体里，不同频率的波速度略有差异（色散），长距离传播后波形会被“拉开”。地震学家分析地震波、医生做超声成像，都必须把这种差异算进去。

第三，多普勒公式 \(f'=vf/(v-u)\) 的适用范围写得明明白白：\(u<v\)。声源达到或超过声速，公式失效，冲击波模型接管。更重要的是，别把它直接搬给光。声音的多普勒效应本质上是“波相对于介质的赛跑”，而光不需要介质——没有那条可以赛跑的跑道。光的多普勒效应要用相对论重新推导，结果与声音公式在低速时相近、高速时根本不同。这是后面相对论章节的故事。

第四，分贝不是纯物理量，而是物理量（强度）经过人耳感知规律折算后的刻度。而且人耳对不同频率的灵敏度差别很大：同样强度，三千赫兹听起来远比五十赫兹响。所以“响度”一半是物理学，一半是生理学。

第五，“声音需要介质”这句话在介质稀薄到一定程度时，不是“声音变弱”，而是“波这个图像整个垮掉”。当空气稀薄到分子平均要飞过比波长还远的距离才能撞上另一个分子时，“挤一下传给下一个”的接力根本组织不起来，也就没有声波可言。高层大气和星际空间都属于这种情况。

最后是本节的典型误用清单。“嗓门大所以传得快”——错，振幅不影响声速，大声只是传得“远”（能量多，衰减到听不见需要更长的路）。“音调高所以刺耳费能量”——错，音调和响度是两个独立的量。“真空中声音很小”——错，真空中没有声音，这不是程度问题，是存在与否的问题。“顺风时风把声音吹了过去”——这个画面也不对：风的速度比声速小得多，根本驮不动声音；真正的原因是高空的风通常比地面快，声波的上半部分被推着多走了一点，整个波前被“压弯”回向地面，声音因此传得更远。这是折射，不是快递。
\end{mainline}

\step{回头解释开场现象}

\begin{mainline}
现在回到那片可怜的扬声器膜。

谜底的第一步：空气里的叠加。乐队几十件乐器同时在空气里掀起各自的疏密波，但在你耳朵所在的那个点上，压强在每一时刻只有一个数值——所有乐器各自造成的压强起伏，在那个点上加成一条随时间起伏的总曲线。整首交响乐，对空气而言只是一条曲线：压强随时间的一张图。这就是为什么一片膜够用——它不需要“同时做几十件事”，它只需要在每一时刻把自己摆在正确的位置，把那条总曲线一丝不差地复演一遍。第三注的答案揭晓：录音存的不是乐器名单，而是这条总曲线每间隔极短时间的一个数值（常见的是每秒四万四千一百个）。播放时把数值序列变回膜的位置序列，整条曲线原样复活。

谜底的第二步在你脑袋里。既然到达鼓膜的只有一条缠成一团的曲线，你怎么还能把鼓、吉他和人声一样一样分出来？因为你的内耳里有一块叫耳蜗的结构，里面有一张基底膜：它窄而硬的一端对高频共振，宽而软的一端对低频共振。声波进来后，不同频率的成分在基底膜的不同位置激起振动，各自触动不同的神经——物理上完成了一次“频谱分析”。它像一台频谱分析仪，像的地方在于：它确实把叠加的曲线按频率拆开、分行上报。不像的地方在于：它的分辨率有限，相近的频率会互相“盖住”（掩蔽），而且上报之后，大脑还要根据经验把分行信息重新拼成“这是小提琴”“这是他的声音”。音乐压缩格式（如 MP3）正是钻掩蔽效应的空子：把那些注定被强音盖住、你本来也听不到的频谱成分直接删掉，文件缩到十分之一，你几乎听不出差别。这是你猜不到的第三层：连“听清”本身，都是一场物理分析和生理猜谜的接力。

于是“信息”这个概念浮出水面。声音之所以能告诉你“是谁、说了什么、心情如何”，是因为那条压强曲线可以在无数种可能的形状里取其中一种——形状的变化余地越大，能区分的消息就越多，携带的信息就越多。声带和舌头负责把意思“写”进波形，空气负责搬运，耳蜗和大脑负责“读”。写、传、读，就是一切通信的三个环节。回头看第一注和第二注：玻璃罩抽走空气，抽掉的是“传”这一环，铃声自然消失；大声和小声走的是同一条路、同一种接力，速度一模一样，分不出先后——第二注里押“一样长”的人赢了。

这个视角立刻解释了为什么本章的技术应用多得不计其数。电话是把波形变成电流再变回来；语音识别是让计算机代替大脑做“读”这一步；倒车雷达和声呐是主动发出声音、用回声的时间差测量距离——蝙蝠在黑暗里做的正是这件事；医学超声用听不见的高频声波给身体内部“拍照”；主动降噪耳机则把叠加原理反过来用：测出外界噪声的波形，实时发出一条处处相反的波形，两条曲线相加，噪声被“减去”。音乐厅的设计师玩的是另一种平衡：他们既要墙壁把声音反射给后排听众，又不能让回声拖得太长，把前一个音和后一个音糊成一团——于是墙面、座椅、天花板的角度和材料，全是在和回声的时间差讨价还价。

而最深远的后果，藏在本章最大的限制里。声音需要介质，所以声音永远出不了大气层。月球上的宇航员面对面站着，喊破喉咙对方也听不见一丝一毫——他们通话用的是无线电。无线电波和声音一样能携带波形、能调制出信息，却不需要任何介质。不需要介质的“波”，到底是什么东西在“抖”？追问这个问题，会把物理学从机械振动推向一个全新的概念——场。那是下一篇电与磁的起点，也是手机、广播、Wi-Fi 和光本身共同的起点。
\end{mainline}

\step{脑洞实验与挑战题}

\begin{thought}[如果星际空间充满空气]
太阳表面时刻都在剧烈地震动，天文学家甚至能通过观测太阳表面的起伏来“听诊”太阳内部，这门学问叫星震学。现在开个脑洞：假如太阳和地球之间的太空充满了和地面一样稠密的空气，会发生什么？

先算时间。光从太阳到地球约八分钟；声音每秒三百四十米，走完同样的一亿五千万公里要约一万四千年。也就是说，如果那位神话里的传令官靠嗓子报告“太阳出来了”，这条消息现在还在半路上，要再等一万多年。

再算响度。声音强度随距离平方衰减，而日地距离是天文数字——就算太阳喊得再响，抵达地球时也会衰减到听不见的程度。更何况真实的星际空间接近真空，连这条衰减的路都不存在。

所以这个脑洞的结论是双重的。坏消息：宇宙的“声音”被永久静音了，再壮观的恒星爆发也传不来一丝声响。好消息：振动本身仍然可以被“看见”——太阳表面的起伏会改变它发出的光，我们用光当信使，间接读出了太阳的“声音”。声音到不了的地方，光能到。为什么光能到？这个问题价值连城，它的答案就是下一篇。
\end{thought}

\begin{puzzles}
\item 雷雨天，你看到闪电后数到六秒才听到雷声，于是判断雷电在两公里外。朋友质疑：“万一声速不是每秒三百四十米呢？”请回答：声速在空气里会不会变化？受什么影响？在什么情况下你的“两公里”会明显估错？（思路提示：回忆声速由介质的什么性质决定；想想冬天和夏天、山脚下和高山上空气有什么差别。光的飞行时间为什么可以放心忽略？）

\item 棉线电话正通着话，第三个人用手指在棉线中点轻轻捏了一下，通话为什么立刻中断？如果在棉线中点再牢牢接上第三个纸杯，做成“三方通话”，你觉得能成吗？音量会怎么变？（思路提示：振动传到哪里，能量就跟到哪里；被捏住的点振动还自由吗？多接一条路，原来的能量要分成几份？）

\item 一辆救护车以每秒三十米（约时速一百一十公里）行驶，警笛频率五百赫兹。取声速每秒三百四十米，算出它朝你开来和离你远去时你听到的频率。然后用两个特殊检查验证你的代数值是否合理：其一，车速为零时应回到五百赫兹；其二，接近时的频率应该高于远离时。（思路提示：直接代入 \(f'=vf/(v-u)\)，远离时 \(u\) 取负；检查不是形式步骤，每个检查都要说出“如果不满足就说明哪里错了”。）

\item 用锤子敲一下桌面，是一声没有明确音调的“哒”；敲一下音叉，是音调很准的长音。两种声音都是物体振动产生的，为什么一个有音调、一个没有？（思路提示：音调对应频率，而一段波形的“频率明确程度”和它的持续时间有关——想想挑战层里那条 \(\Delta f\,\Delta t\gtrsim 1\)。再想想：如果这个规律对所有波都成立，它对后面要学的微观世界的波意味着什么？）
\end{puzzles}
